\chapter{电动力学预备班}
\epigraph{你应该已经学过微积分和矢量分析了, 对吧?}{——ZexHexoft}
这部分主讲常用矢量分析和一些相关的定理及公式.\footnote{事实上我这里非常推荐阅读《电磁学(拓展篇)》中的专题15, 接下来也会引用部分内容}

\section{常用的矢量恒等式}

\begin{theorem}[标量三重积]
    \begin{equation}
        \vec A\cdot(\vec B\times\vec C)=\vec B\cdot(\vec C\times\vec A)=\vec C\cdot(\vec A\times\vec B)
    \end{equation}
\end{theorem}

\begin{theorem}[矢量三重积(BAC-CAB规则\footnote{英文中与BACK-CAB同音, 即后面的出租车})] \label{BAC-CAB}
        \begin{equation}
            \vec A\times(\vec B\times\vec C)=\vec B(\vec A\cdot\vec C)-\vec C(\vec A\cdot\vec B)
        \end{equation}
\end{theorem}

\section{矢量的微分}

\subsection{一阶微分}

\begin{definition}[$\nabla$算子]
    \begin{equation}
        \nabla := \hat{\mathbf{x}}\frac{\partial}{\partial x}+\hat{\mathbf{y}}\frac{\partial}{\partial y}+\hat{\mathbf{z}}\frac{\partial}{\partial z}
    \end{equation}
\end{definition}

\begin{theorem}[梯度]
    \begin{equation}
        \nabla T := \frac{\partial T}{\partial x}\hat{\mathbf{x}}+\frac{\partial T}{\partial y}\hat{\mathbf{y}}+\frac{\partial T}{\partial z}\hat{\mathbf{z}}
    \end{equation}
\end{theorem}

\begin{theorem}[散度]
    \begin{equation}
        \nabla \cdot \mathbf{v}=\frac{\partial v_x}{\partial x}+\frac{\partial v_y}{\partial y}+\frac{\partial v_z}{\partial z}
    \end{equation}
\end{theorem}

在\cref{sec:integrationOfVector}中，我们会看到他们名字的意义.

\begin{theorem}[旋度]\footnote{都学到这里了，你会计算这个的，对吧? 那我就不展开了.}
    \begin{equation}
        \nabla \times \mathbf{v}=\left| \begin{matrix}
            \hat{\mathbf{x}}&		\hat{\mathbf{y}}&		\hat{\mathbf{z}}\\
            \frac{\partial}{\partial x}&		\frac{\partial}{\partial y}&		\frac{\partial}{\partial z}\\
            v_x&		v_y&		v_z\\
        \end{matrix} \right|
    \end{equation}
\end{theorem}

\begin{theorem}[微分的积规则]
    \begin{equation}
        \nabla(fg) = f\nabla g+g\nabla f
    \end{equation}
    \begin{equation}
        \nabla(\vec A \cdot \vec B) = \vec A \times (\nabla \times \vec B) + \vec B \times (\nabla \times \vec A) + (\vec A \cdot \nabla)\vec B+ (\vec B \cdot \nabla)\vec A \label{eq:example2}
    \end{equation}
    \begin{equation}
        \nabla \cdot \left( f\mathbf{A} \right) =f\left( \nabla \cdot \mathbf{A} \right) +\mathbf{A}\cdot \nabla f
    \end{equation}
    \begin{equation}
        \nabla \cdot \left( \mathbf{A}\times \mathbf{B} \right) =\mathbf{B}\cdot \left( \nabla \times \mathbf{A} \right) -\mathbf{A}\cdot \left( \nabla \times \mathbf{B} \right) 
    \end{equation}
    \begin{equation}
        \nabla \times \left( f\mathbf{A} \right) =f\left( \nabla \times \mathbf{A} \right) -\mathbf{A}\times \left( \nabla f \right) 
    \end{equation}
    \begin{equation}
        \nabla \times \left( \mathbf{A}\times \mathbf{B} \right) =\left( \mathbf{B}\cdot \nabla \right) \mathbf{A}-\left( \mathbf{A}\cdot \nabla \right) \mathbf{B}+\mathbf{A}\left( \nabla \cdot \mathbf{B} \right) -\mathbf{B}\left( \nabla \cdot \mathbf{A} \right) \label{eq:example1}
    \end{equation}
\end{theorem}

说实话，这些等式又臭又长，很多时候难以记住.

我在这里提供一种摘自《电磁学(拓展篇)》中的专题15中选读15-10的一种记忆方法供大家参考.

首先简记一种符号$\nabla_\vec{A}$(可以与一元微积分中的$\frac{\mathrm{d}f}{\mathrm{d}x}\frac{\partial}{\partial f},\frac{\mathrm{d}g}{\mathrm{d}x}\frac{\partial}{\partial g}$等对应)

\begin{definition}
    $\nabla_\vec{A}$表示只作用于向量$\vec{A}$的算符.
\end{definition}

这里以\cref{eq:example1}为例子

\begin{example}
    \begin{align*}
        \nabla \times \left( \mathbf{A}\times \mathbf{B} \right) &=\nabla _{\mathbf{A}}\times \left( \mathbf{A}\times \mathbf{B} \right) +\nabla _{\mathbf{B}}\times \left( \mathbf{A}\times \mathbf{B} \right) 
        \\
        &=\mathbf{A}\left( \nabla _{\mathbf{A}}\cdot \mathbf{B} \right) -\mathbf{B}\left( \nabla _{\mathbf{A}}\cdot \mathbf{A} \right) +\mathbf{A}\left( \nabla _{\mathbf{B}}\cdot \mathbf{B} \right) -\mathbf{B}\left( \nabla _{\mathbf{B}}\cdot \mathbf{A} \right)
        \\
        &=\left( \mathbf{B}\cdot \nabla _{\mathbf{A}} \right) \mathbf{A}-\mathbf{B}\left( \nabla _{\mathbf{A}}\cdot \mathbf{A} \right) +\mathbf{A}\left( \nabla _{\mathbf{B}}\cdot \mathbf{B} \right) -\left( \mathbf{A}\cdot \nabla _{\mathbf{B}} \right) \mathbf{B}
        \\
        &=\left( \mathbf{B}\cdot \nabla \right) \mathbf{A}-\mathbf{B}\left( \nabla \cdot \mathbf{A} \right) +\mathbf{A}\left( \nabla \cdot \mathbf{B} \right) -\left( \mathbf{A}\cdot \nabla \right) \mathbf{B}
    \end{align*}
\end{example}

再举一个\cref{eq:example2}的例子(注意对\cref{BAC-CAB}的逆运用)

\begin{example}
    \begin{align*}
        \nabla \left( \mathbf{A}\cdot \mathbf{B} \right) &=\nabla _{\mathbf{A}}\left( \mathbf{A}\cdot \mathbf{B} \right) +\nabla _{\mathbf{B}}\left( \mathbf{A}\cdot \mathbf{B} \right)
        \\
        &=\mathbf{B}\times \left( \nabla _{\mathbf{A}}\times \mathbf{A} \right) +\mathbf{A}\left( \mathbf{B}\cdot \nabla _{\mathbf{A}} \right) +\mathbf{A}\times \left( \nabla _{\mathbf{B}}\times \mathbf{B} \right) +\mathbf{B}\left( \mathbf{A}\cdot \nabla _{\mathbf{B}} \right) 
        \\
        &=\mathbf{B}\times \left( \nabla \times \mathbf{A} \right) +\left( \mathbf{B}\cdot \nabla \right) \mathbf{A}+\mathbf{A}\times \left( \nabla \times \mathbf{B} \right) +\left( \mathbf{A}\cdot \nabla \right) \mathbf{B}
    \end{align*}
\end{example}

\subsection{常用的二阶微分}

\begin{definition}[拉普拉斯算符]\footnote{似乎有点反直觉的是，拉普拉斯算符既可以作用在标量上，也可以作用在矢量上.}
    \begin{equation}
        \nabla^2 := \nabla \cdot \nabla = \frac{\partial ^2}{\partial x^2}+\frac{\partial ^2}{\partial y^2}+\frac{\partial ^2}{\partial z^2}
    \end{equation}
\end{definition}

\begin{theorem}
    \begin{gather}
        \nabla \times \left( \nabla T \right) =0\\
        \nabla \cdot \left( \nabla \times \mathbf{v} \right) =0\\
        \nabla \times \left( \nabla \times \mathbf{v} \right) =\nabla \left( \nabla \cdot \mathbf{v} \right) -\nabla ^2\mathbf{v}
    \end{gather}
\end{theorem}


\section{矢量的积分}\label{sec:integrationOfVector}

\begin{theorem}[梯度定理]
    \begin{equation}
        \int_a^b{\left( \nabla T \right) \cdot \mathrm{d}\mathbf{l}}=T\left( b \right) -T\left( a \right) 
    \end{equation}
\end{theorem}

\begin{theorem}[散度定理]
    \begin{equation}
        \int_V{\left( \nabla \cdot \mathbf{v} \right) \mathrm{d}\tau}=\oint_S{\mathbf{v}\cdot \mathrm{d}\mathbf{a}}
    \end{equation}
\end{theorem}

\begin{theorem}[旋度定理/斯托克斯定理]\footnote{这里的$\partial S$表示$S$的边界.}\footnote{之所以在积分号上用逆时针的圆圈，是为了强调右手规则.}
    \begin{equation}
        \int_S{\left( \nabla \times \mathbf{v} \right) \cdot \mathrm{d}\mathbf{a}}=\ointctrclockwise_{\partial S}{\mathbf{v}\cdot \mathrm{d}\mathbf{l}}
    \end{equation}
\end{theorem}